\section{From a weight on a scale to a recovery percentage}

Two different equations in this project both call themselves the recovery model, and they
look nothing alike. Before anything else, here is which is which.

The equation this report is about is fitted in \emph{real units}, with current density $J$ in
$\mathrm{mA/cm^2}$ and the feed ratio $R$ in $\mathrm{mol/mol}$, and it keeps all five terms
that these runs can settle:

\begin{equation}
  \text{recovery \%} = 33.0192 + 3.16999\,J - 14.4538\,R - 0.0954576\,J^{2} + 0.837907\,R^{2}
  \label{eq:real-units-one-run}
\end{equation}

The equation \texttt{uv run limodel analyze} prints is not that one. The code logs both
settings, codes them onto $-1\ldots1$, and then chooses its terms by forward stepwise
selection, so it drops every term that does not lower the leave-one-out error. What it prints
is

\begin{equation}
  \text{recovery \%} = 11.5776 - 15.627\,x_{2} + 6.08016\,x_{2}^{2},
  \qquad x_{2} = \frac{\ln R - 1.395778}{1.039721}
  \label{eq:coded-units-fit}
\end{equation}

Different units, different terms, and no current density at all in \eqref{eq:coded-units-fit}.
The next section sets them side by side. What matters here is that
\eqref{eq:real-units-one-run} and \eqref{eq:coded-units-fit} are fitted to \emph{the same nine
numbers} on the left-hand side: nine measured recovery percentages, one per GO-10 run. Neither
equation can be better than those nine numbers, so this section takes one of them apart
completely. Run AE's recovery is $38.27\%$. Below is every reading that figure is made of and
every multiplication done to it, so it can be checked on paper.

\subsection{The rows AE's recovery is built from}

Nothing in the data files is a recovery percentage. Four raw rows across three CSV files hold
everything \eqref{eq:real-units-one-run} is fitted to for run AE, and Table~\ref{tab:ae-raw-rows}
is all of them. The tares come from the run sheet, the weights from the balance under each
tank, the concentrations from the lab.

\begin{table}[htbp]
\centering
\caption{The raw rows behind AE's recovery, as they sit in the data files. The two tares are
columns of AE's own row in \texttt{runs.csv}; each tank weight is a row of
\texttt{tank\_readings.csv}; each concentration is a row of \texttt{samples.csv}.}
\label{tab:ae-raw-rows}
\begin{tabular}{llrr}
\toprule
file & column or row & minute 0 & minute 240 \\
\midrule
\texttt{runs.csv}          & \texttt{tare\_conc\_g}  & \multicolumn{2}{c}{$96.2\unit{g}$} \\
                           & \texttt{tare\_dil\_g}   & \multicolumn{2}{c}{$111.5\unit{g}$} \\
\addlinespace
\texttt{tank\_readings.csv}& concentrate \texttt{weight\_g} & $2119.7\unit{g}$ & $2101.5\unit{g}$ \\
                           & dilute \texttt{weight\_g}      & $2117.8\unit{g}$ & $2040.7\unit{g}$ \\
\addlinespace
\texttt{samples.csv}       & concentrate \texttt{li\_mg\_l}
                           & $63.56\unit{mg/L}$ & $218.5\unit{mg/L}$ \\
                           & dilute \texttt{li\_mg\_l}
                           & $403.1\unit{mg/L}$ & $261.9\unit{mg/L}$ \\
\bottomrule
\end{tabular}
\end{table}

Two of the eight entries never reach recovery. The dilute's minute 240 weight and lithium are
in the table only to be set aside: recovery divides by what the feed held \emph{before}
anything moved, so only the dilute's first sampled minute is used. Its last minute is read by
the balance gap, which is a diagnostic and is never fitted.

\subsection{Weights become litres}

No instrument measured a volume. Each tank sat on a balance for the whole run, so the litres
are inferred: the balance reading less that tank's own empty weight, divided by
$1000\unit{g/L}$. \texttt{VolumeInATank.find\_litres} does this, and the divisor is
\texttt{SOLUTION\_\allowbreak G\_\allowbreak PER\_\allowbreak L} $= 1000.0$ in \texttt{calculate/\allowbreak constants.py}. That constant is a
decision, not a measurement. These are dilute salt solutions, so a litre weighs a kilogram
closely enough that the balance can be read as a volume, and the alternative --- a density per
tank per minute --- was never measured.

The concentrate is weighed at both ends, so it needs two volumes:

\begin{align}
  V_{\mathrm{conc}}(0)   &= \frac{2119.7 - 96.2}{1000} = 2.0235\unit{L} \label{eq:v-conc-0} \\
  V_{\mathrm{conc}}(240) &= \frac{2101.5 - 96.2}{1000} = 2.0053\unit{L} \label{eq:v-conc-240}
\end{align}

The tare is looked up by tank name through \texttt{TARE\_FIELD\_BY\_TANK}, which maps
\texttt{conc} to the column \texttt{tare\_conc\_g}; \texttt{AmountThatCrossed.find\_ends} is
what reads it off the run. The two volumes differ by $18.2\unit{g}$ of liquid, nine parts in a
thousand. That is small, and it is exactly why the volume is recomputed at each sampled minute
instead of once: a concentration at minute 240 compared against a minute 0 volume would be
comparing two different quantities of liquid.

\subsection{Litres and a lab result become milligrams}

A concentration on its own says nothing that can be compared across a run, because the tanks
gain and lose liquid all the way through. Only concentration times volume --- an absolute
amount --- can be subtracted from one minute to the next.
\texttt{AmountInATank.find\_milligrams} does the multiplication, taking the lithium column
name from \texttt{LITHIUM.field\_name}, which is \texttt{li\_mg\_l}:

\begin{align}
  m_{\mathrm{conc}}(0)   &= 63.56 \times 2.0235 = 128.61366\unit{mg} \label{eq:mg-conc-0} \\
  m_{\mathrm{conc}}(240) &= 218.5 \times 2.0053 = 438.15805\unit{mg} \label{eq:mg-conc-240}
\end{align}

The concentrate did not start empty of lithium. It was made up as $0.5\unit{g/L}$ LiCl, so
\eqref{eq:mg-conc-0} is not a rounding artefact and cannot be skipped --- about $29\%$ of what
the concentrate finished with was already in it. What crossed the membranes is the difference,
and \texttt{AmountThatCrossed.\allowbreak find\_milligrams} is the subtraction:

\begin{equation}
  \Delta m_{\mathrm{conc}} = 438.15805 - 128.61366 = 309.54439\unit{mg}
  \label{eq:crossed}
\end{equation}

The two ends are picked by minute and not by row order.
In \texttt{AmountThatCrossed}, the method
\texttt{find\_amounts\_by\_minute} builds a dictionary keyed on the minute, keeping only
the minutes that have both a weight and a lab result, and \texttt{find\_ends} then sorts
those keys and takes the first and the last. The readings arrive in no guaranteed order, and
most minutes of AE have a weight but no sample, so taking the first and last \emph{rows}
would pick up the wrong pair.

\subsection{The feed side, and the percentage}

The denominator is the whole prize: all the lithium the feed held when the run began. It is
built the same way, from the dilute tank's own tare and its first sampled minute:

\begin{align}
  V_{\mathrm{dil}}(0) &= \frac{2117.8 - 111.5}{1000} = 2.0063\unit{L} \label{eq:v-dil-0} \\
  m_{\mathrm{dil}}(0) &= 403.1 \times 2.0063 = 808.73953\unit{mg} \label{eq:mg-dil-0}
\end{align}

\texttt{LithiumRecovery.find\_percent} then divides \eqref{eq:crossed} by
\eqref{eq:mg-dil-0}:

\begin{equation}
  \text{recovery \%} = 100 \times \frac{309.54439}{808.73953} = 38.2749\%
  \label{eq:recovery}
\end{equation}

Three things about \eqref{eq:recovery} were choices. First, the denominator is the dilute's
\emph{first} sampled minute and never its last: measured against what the feed ended with,
AE would read $61.3\%$, which would be a share of what was left rather than a recovery.
Second, the numerator is what the \emph{concentrate gained}, not what the dilute lost --- the
concentrate is where the product is, and the dilute also loses lithium to the rinse and to the
sample bottles. Those two numbers disagree by a measurable amount, and the gap between them is
reported separately as a diagnostic rather than averaged away. Third, milligrams stand on both
sides of the division, so lithium's molar mass cancels and the code never converts to moles
here; recovery is the one quantity in the project that needs no molar mass at all.

Figure~\ref{fig:ae-recovery-chain} is the whole chain with these numbers on it.

\begin{figure}[htbp]
\centering
\footnotesize
\begin{tikzpicture}[
  every node/.style={font=\footnotesize},
  box/.style={draw, rounded corners, align=center, text width=2.0cm, inner sep=3pt,
              minimum height=1.0cm},
  lbl/.style={font=\scriptsize, align=center, inner sep=1pt},
  arr/.style={-Stealth, thick}
]

\node[box] (cw0) at (0,2.5) {concentrate\\min 0\\$2119.7\unit{g}$};
\node[box] (cw4) at (0,0.4) {concentrate\\min 240\\$2101.5\unit{g}$};
\node[box] (dw0) at (0,-2.3) {dilute\\min 0\\$2117.8\unit{g}$};

\node[box] (cv0) at (4.2,2.5) {$2.0235\unit{L}$};
\node[box] (cv4) at (4.2,0.4) {$2.0053\unit{L}$};
\node[box] (dv0) at (4.2,-2.3) {$2.0063\unit{L}$};

\node[box] (cm0) at (8.4,2.5) {$128.61\unit{mg}$\\lithium};
\node[box] (cm4) at (8.4,0.4) {$438.16\unit{mg}$\\lithium};
\node[box] (dm0) at (8.4,-2.3) {$808.74\unit{mg}$\\lithium};

\node[box, text width=2.3cm] (cross) at (12.7,1.45)
  {crossed\\$438.16 - 128.61$\\$= 309.54\unit{mg}$};
\node[box, text width=2.2cm, very thick] (rec) at (12.7,-2.3) {recovery\\$38.27\%$};

\draw[arr] (cw0) -- node[lbl, above] {less $96.2$\\$\div\,1000$} (cv0);
\draw[arr] (cw4) -- node[lbl, above] {less $96.2$\\$\div\,1000$} (cv4);
\draw[arr] (dw0) -- node[lbl, above] {less $111.5$\\$\div\,1000$} (dv0);

\draw[arr] (cv0) -- node[lbl, above] {$\times\,63.56$} (cm0);
\draw[arr] (cv4) -- node[lbl, above] {$\times\,218.5$} (cm4);
\draw[arr] (dv0) -- node[lbl, above] {$\times\,403.1$} (dm0);

\draw[arr] (cm4) -- (cross);
\draw[arr] (cm0) -- (cross);
\draw[arr] (cross) -- node[lbl, right] {$\div\,808.74$\\$\times\,100$} (rec);
\draw[arr] (dm0) -- (rec);

\node[lbl, anchor=west] at (-1.35,3.7) {\textbf{balance}};
\node[lbl, anchor=west] at (2.85,3.7) {\textbf{litres}};
\node[lbl, anchor=west] at (7.05,3.7) {\textbf{milligrams}};
\node[lbl, anchor=west] at (11.3,3.7) {\textbf{outcome}};

\end{tikzpicture}
\caption{Run AE's recovery, every step with its own number. The tares come from
\texttt{runs.csv}, the weights from \texttt{tank\_readings.csv}, the multipliers from the
lithium column of \texttt{samples.csv}. Left to right the classes are
\texttt{VolumeInATank}, \texttt{AmountInATank}, \texttt{AmountThatCrossed} and
\texttt{LithiumRecovery}.}
\label{fig:ae-recovery-chain}
\end{figure}

\subsection{Permselectivity, which needs magnesium and the charge}

Recovery only asks how much lithium moved. Permselectivity asks whether the stack
\emph{preferred} lithium, which takes magnesium as well, and it has one wrinkle recovery does
not: it is about charge, not about ions.

The magnesium that crossed comes through exactly the same chain. The only change is the
column read, \texttt{MAGNESIUM.field\_name}, which is \texttt{mg\_mg\_l}:

\begin{equation}
  \Delta m_{\mathrm{Mg}} = 130.0 \times 2.0053 \;-\; 2.974 \times 2.0235
  = 260.689 - 6.01789 = 254.67111\unit{mg}
  \label{eq:crossed-mg}
\end{equation}

\texttt{TransportNumbers.find\_numbers} then converts both to millimoles and weights each by
the charge its own ion carries --- \texttt{Metal.charge}, which is $1$ for lithium and $2$ for
magnesium. Magnesium counts twice because a transport number is a share of the charge the
current carried, and one $\mathrm{Mg}^{2+}$ ion carries two charges across the membrane where
one $\mathrm{Li}^{+}$ carries one. The current is what was paid for, so the current is what
gets divided up:

\begin{align}
  q_{\mathrm{Li}} &= 1 \times \frac{309.54439}{6.94}   = 44.60294 \label{eq:q-li} \\
  q_{\mathrm{Mg}} &= 2 \times \frac{254.67111}{24.305} = 20.95627 \label{eq:q-mg} \\
  t_{\mathrm{Li}} &= \frac{44.60294}{65.55921} = 0.680346,
  \qquad t_{\mathrm{Mg}} = \frac{20.95627}{65.55921} = 0.319654 \label{eq:transport}
\end{align}

So lithium carried $68\%$ of the charge into the concentrate. On its own that number says
little, because AE's feed held more lithium than magnesium in molar terms and a stack with no
preference at all would still move mostly lithium. The feed it was offered has to be divided
out. \texttt{FeedRatio.find\_li\_over\_mg\_in\_the\_dilute} takes the first dilute minute the
lab reported both metals for, in millimoles. That minute is AE's minute 0, where the dilute
held $808.73953\unit{mg}$ of lithium by \eqref{eq:mg-dil-0} and
$1657 \times 2.0063 = 3324.4391\unit{mg}$ of magnesium:

\begin{equation}
  \frac{c_{\mathrm{Li}}}{c_{\mathrm{Mg}}}
  = \frac{808.73953 / 6.94}{3324.4391 / 24.305}
  = \frac{116.53307}{136.78005} = 0.851974
  \label{eq:feed-ratio}
\end{equation}

Both amounts are from one tank at one minute, so the litres cancel and the milligrams carry
the ratio --- which is why \eqref{eq:feed-ratio} divides milligrams-over-molar-mass directly.
\texttt{Permselectivity.\allowbreak find\_value} finishes it:

\begin{equation}
  P = \frac{t_{\mathrm{Li}} / t_{\mathrm{Mg}}}{c_{\mathrm{Li}} / c_{\mathrm{Mg}}}
    = \frac{0.680346 / 0.319654}{0.851974}
    = \frac{2.128381}{0.851974} = 2.498176
  \label{eq:permselectivity}
\end{equation}

\subsubsection*{Why no preference is $0.5$ and not $1$}

Dividing by the feed ratio is what lets runs fed different mixtures be compared, and it fixes
the neutral point at $0.5$ whatever the feed was. The reason is the factor of two in
\eqref{eq:q-mg}, which survives the division.

Suppose the stack had no preference, so the ions crossed in proportion to the moles offered:
$n_{\mathrm{Li}} / n_{\mathrm{Mg}} = c_{\mathrm{Li}} / c_{\mathrm{Mg}}$. The transport numbers
are charge shares, so

\begin{equation}
  \frac{t_{\mathrm{Li}}}{t_{\mathrm{Mg}}}
  = \frac{1 \cdot n_{\mathrm{Li}}}{2 \cdot n_{\mathrm{Mg}}}
  = \frac{1}{2}\,\frac{c_{\mathrm{Li}}}{c_{\mathrm{Mg}}},
  \qquad\text{so}\qquad
  P = \frac{\tfrac{1}{2}\,(c_{\mathrm{Li}} / c_{\mathrm{Mg}})}{c_{\mathrm{Li}} / c_{\mathrm{Mg}}}
    = \frac{1}{2}
  \label{eq:neutral-point}
\end{equation}

The feed ratio cancels completely, so $0.5$ is the neutral point for every run regardless of
what it was fed. For AE, no preference would have meant
$t_{\mathrm{Li}} / t_{\mathrm{Mg}} = 0.851974 / 2 = 0.425987$; the stack actually reached
$2.128381$, about five times that, hence $P = 2.498$ against a neutral $0.5$. Read against $1$
instead, the same stack would look only $2.5$ times selective rather than $5$, so the neutral
point has to be stated with the number.

One more thing about \eqref{eq:feed-ratio}: it is the \emph{measured} feed ratio, from the lab
result on the dilute sample. It is not the ratio the run was set up at. AE was made up at
$400\unit{mg/L}$ lithium and $2000\unit{mg/L}$ magnesium, which is
$(2000/24.305) / (400/6.94) = 1.427690$ mol magnesium per mole lithium, or $0.700432$ the
other way up --- not $0.851974$. \texttt{FeedRatio} keeps the two apart by name,
\texttt{find\_li\_over\_mg\_in\_the\_dilute} against
\texttt{find\_mg\_over\_li\_from\_the\_setup}, because they answer different questions. The
measured one is what permselectivity divides by, since it is the mixture the membrane actually
saw. The set-up one is the design column $R$ of \eqref{eq:real-units-one-run}, since a design
column has to be a setting the bench chose and can choose again, not something a lab result
moves afterwards.

\subsection{What else a \texttt{Result} carries}

\texttt{Calculator.find\_results} builds one \texttt{Result} per run, and it drops a run only
when permselectivity or recovery could not be worked out --- those two are what the model is
fitted to, so a run without them has nothing to contribute. The rest may be \texttt{None} on a
run that is otherwise sound, which is why they are reported at the end rather than used to
drop anything. Table~\ref{tab:result-fields} names each field, the class behind it, and what
AE came out as.

\begin{table}[htbp]
\centering
\caption{Every field of a \texttt{Result} beyond the two derived above, with the class that
produces it and AE's value. The last five are diagnostics: read by a human judging a run,
never fitted.}
\label{tab:result-fields}
\begin{tabular}{llr}
\toprule
field & class and method & AE \\
\midrule
\texttt{current\_density\_ma\_cm2} & \texttt{CurrentDensity.find\_ma\_per\_cm2} & $13.6667$ \\
\texttt{feed\_mg\_li\_ratio}
  & \texttt{FeedRatio.find\_mg\_over\_li\_from\_the\_setup} & $1.42769$ \\
\addlinespace
\texttt{highest\_voltage\_v} & \texttt{StackVoltage.find\_highest\_volts} & $16.0\unit{V}$ \\
\texttt{charge\_passed\_mah\_cm2} & \texttt{ChargeThatPassed.find\_mah\_per\_cm2} & $54.5625$ \\
\texttt{energy\_per\_mol\_li\_wh} & \texttt{EnergyPerMoleOfLithium.find\_watt\_hours} & $91.92$ \\
\texttt{li\_balance\_gap\_pct} & \texttt{BalanceGap.find\_percent} & $-19.64\%$ \\
\texttt{mg\_balance\_gap\_pct} & \texttt{BalanceGap.find\_percent} & $35.00\%$ \\
\bottomrule
\end{tabular}
\end{table}

Four of these are worth a line on why they are shaped the way they are.

\textbf{Current density} is the \emph{median} supply current over the power log, times $1000$
for milliamps, divided by the membrane area. AE's power log carries a current at 30
minutes, 25 of them $0.082\unit{A}$, so $0.082 \times 1000 / 6 = 13.6667\unit{mA/cm^2}$.
The median and not the
mean, because a run holds one current all the way through and a single misread minute should
not shift the number that labels the whole run. This is the $J$ of
\eqref{eq:real-units-one-run}.

\textbf{Highest voltage} is the largest voltage the log ever recorded, which is what the
$20\unit{V}$ ceiling is checked against. \texttt{StackVoltage} also offers
\texttt{find\_starting\_volts}, the voltage at the earliest logged minute, which is what the
stack asked for before anything had crossed.

\textbf{Charge passed} and \textbf{energy per mole} are both areas under the power log, found
by the trapezium rule between consecutive logged minutes, because no meter totalled either
one. AE integrates to $1178.55\unit{C}$, which over $6\unit{cm^2}$ is
$54.5625\unit{mAh/cm^2}$, and to $14759.78\unit{J}$, which against the
$309.54\unit{mg} / 6.94 = 44.6\unit{mmol}$ of lithium that crossed is
$91.92\unit{Wh}$ per mole. Per mole of lithium rather than per hour or per litre, so a short
run and a long one can be set beside each other.

\textbf{The balance gaps} are the honest confession of the chain above. Each one is the
percent of what the dilute lost that was not found again in the concentrate plus the rinse.
AE's lithium gap is $-19.64\%$: more lithium was found than the dilute appears to have lost,
which is a measurement inconsistency and not a physical one. It is reported and never
corrected for, and no quantity in this project is adjusted by it. When the rinse was never
weighed it counts as zero gained, which can only make the gap look worse --- the honest
direction for a diagnostic to be wrong in.

With all of that, AE reduces to one row the fitting can see: $J = 13.6667$, $R = 1.42769$,
recovery $= 38.27\%$, permselectivity $= 2.498$. Nine such rows are everything
\eqref{eq:real-units-one-run} and \eqref{eq:coded-units-fit} were fitted to.
